\documentclass[UTF8,11pt,a4paper]{ctexart}
\usepackage[a4paper,left=24mm,right=24mm,top=23mm,bottom=22mm,headheight=15pt]{geometry}
\usepackage{xcolor,booktabs,tabularx,array,enumitem,fancyhdr,hyperref,titlesec,lastpage,tikz,graphicx}
\usetikzlibrary{arrows.meta,positioning}
\graphicspath{{./}{assets/}}
\definecolor{navy}{HTML}{173B57}
\definecolor{blue}{HTML}{247BA0}
\definecolor{cyan}{HTML}{5BC0BE}
\definecolor{ink}{HTML}{23313D}
\definecolor{muted}{HTML}{657786}
\definecolor{paper}{HTML}{F4F7F9}
\hypersetup{colorlinks=true,linkcolor=navy,urlcolor=blue,pdfauthor={Slurm AI Agent 项目组},pdftitle={Slurm AI Agent 作品简介}}
\setlength{\parindent}{2em}
\setlength{\parskip}{0.3em}
\setlength{\emergencystretch}{3em}
\setlist{nosep,leftmargin=2em}
\setlist[itemize]{label=\textcolor{blue}{\small\textbullet}}
\pagestyle{fancy}\fancyhf{}
\fancyhead[L]{\small\textcolor{muted}{Slurm AI Agent 作品简介}}
\fancyhead[R]{\small\textcolor{muted}{USTC 107 算力平台}}
\fancyfoot[C]{\small\textcolor{muted}{第 \thepage\ 页 / 共 \pageref*{LastPage} 页}}
\renewcommand{\headrulewidth}{0.4pt}
\titleformat{\section}{\Large\bfseries\color{navy}}{\thesection}{0.7em}{}[\vspace{-0.2em}\color{cyan}\titlerule]
\titleformat{\subsection}{\large\bfseries\color{blue}}{\thesubsection}{0.6em}{}
\newcommand{\code}[1]{\texttt{\small #1}}

\begin{document}
\begin{titlepage}
\pagecolor{navy}\color{white}
\vspace*{26mm}
{\Huge\bfseries Slurm AI Agent\par}
\vspace{5mm}
{\LARGE 面向 USTC 107 算力平台的智能作业助手\par}
\vspace{13mm}{\color{cyan}\rule{0.72\textwidth}{1.5pt}\par}
\vspace{13mm}{\fontsize{30}{36}\selectfont\bfseries 作品简介\par}
\vspace{9mm}
{\Large 用自然语言降低使用门槛，用确定性后端守住执行边界\par}
\vfill
{\large 项目关键词\par}\vspace{4mm}
\fcolorbox{cyan}{navy}{\strut\quad HPC 智能体\quad}\quad
\fcolorbox{cyan}{navy}{\strut\quad Slurm\quad}\quad
\fcolorbox{cyan}{navy}{\strut\quad Function Calling\quad}\\[3mm]
\fcolorbox{cyan}{navy}{\strut\quad 依赖自愈\quad}\quad
\fcolorbox{cyan}{navy}{\strut\quad 资源核验\quad}\quad
\fcolorbox{cyan}{navy}{\strut\quad 知识检索\quad}
\vspace{22mm}
\end{titlepage}
\nopagecolor\color{ink}

\section{系统概述}
Slurm AI Agent 是部署在 USTC 107 算力平台登录节点上的 Web 工作区：项目文件、Python 环境、作业提交、运行报告，以及接入大模型的智能助手，整合在同一个页面中。用户以自然语言描述任务目标，系统将其转化为真实的 Slurm 与 Conda 操作；每一步都在用户确认与后端校验之下完成。

\begin{figure}[htbp]
\centering
\begin{tikzpicture}[font=\small,
 c/.style={draw=navy,very thick,rounded corners,fill=paper,minimum height=16mm,text width=30mm,align=center},
 s/.style={draw=blue,rounded corners,fill=paper,minimum height=11mm,text width=33mm,align=center},
 l/.style={thick,color=muted}]
 \node[c] (p) at (0,0) {\textbf{项目工作空间}\\部署于登录节点};
 \node[s] (l1) at (-48mm,14mm) {文件管理\\编辑 / 预览 / 上传};
 \node[s] (l2) at (-48mm,0) {项目专属 Conda 环境};
 \node[s] (l3) at (-48mm,-14mm) {智能助手\\19 项受控工具};
 \node[s] (r1) at (48mm,14mm) {依赖检查、确认与安装};
 \node[s] (r2) at (48mm,0) {受控提交与作业监控};
 \node[s] (r3) at (48mm,-14mm) {结果报告与 Web 终端};
 \draw[l] (l1)--(p);\draw[l] (l2)--(p);\draw[l] (l3)--(p);
 \draw[l] (r1)--(p);\draw[l] (r2)--(p);\draw[l] (r3)--(p);
\end{tikzpicture}
\caption{系统全貌：以项目为单元，组织从文件、环境到作业与复盘的全流程。}
\end{figure}

从准备依赖到拿到结果，系统围绕一条闭环展开：

\begin{center}
\begin{tikzpicture}[node distance=5mm and 5mm,font=\small,
 b/.style={draw=blue,rounded corners,fill=paper,minimum height=10mm,text width=31mm,align=center},
 a/.style={-{Latex[length=2mm]},thick,color=muted}]
 \node[b] (p) {项目与依赖准备};
 \node[b,right=of p] (r) {资源规划与受控提交};
 \node[b,right=of r] (m) {排队/运行监控};
 \node[b,right=of m] (o) {结果分析与修正};
 \draw[a] (p)--(r);\draw[a] (r)--(m);\draw[a] (m)--(o);
 \draw[a,bend left=25] (o) to node[above,font=\scriptsize]{失败反馈写回依赖需求} (p);
\end{tikzpicture}
\end{center}

\noindent\begin{tabularx}{\textwidth}{>{\bfseries\color{navy}}p{34mm}X}
\toprule
集群与作业看板 & 展示节点、CPU/GPU 使用、实时作业、个人历史任务和提交选项。\\
项目工作空间 & 创建项目和实验子目录，上传、编辑、预览、下载文件，并维护独立会话。\\
智能助手 & 查询真实集群信息，检索平台文档，读取项目文件，解释配置并分析日志。\\
依赖管理 & 扫描依赖清单、源码 import 和说明文档，生成结构化建议，由用户确认后安装并验证。\\
作业规划与提交 & 提供 PyTorch 单卡/DDP、CPU、Job Array、Jupyter 和通用脚本模板；统一校验后提交。\\
监控与报告 & 持续跟踪作业，结束后生成报告，支持结果文件在线浏览与打包下载。\\
Web 终端与文档 & 在浏览器内使用登录 shell，并阅读或问答 107 平台文档。\\
\bottomrule
\end{tabularx}

% TODO(截图)：此处适合放一张本作品主界面截图（建议：项目工作区总览），放入 latex/assets/ 后取消注释：
% \begin{figure}[htbp]\centering
% \includegraphics[width=\textwidth]{agent-workspace.png}
% \caption{系统主界面：项目工作区与智能助手。}
% \end{figure}

\section{为什么需要它}
107 平台为课程训练、科学计算和深度学习提供了共享算力，但对许多用户而言，从写完代码到拿到结果之间仍有大量环节需要完成：理解分区、账户、QoS 与资源参数的含义，编写 sbatch 脚本，配置 Python 环境，再在作业排队或失败后查找日志、定位原因。一次任务往往需要在平台文档、终端、文件系统和作业队列之间频繁切换，单个环节并不复杂，串联起来却消耗大量时间与注意力。

使用中的高频障碍集中在五类，它们大多并非源于代码本身，而是作业生命周期中繁琐环节过多：
\begin{itemize}
  \item 资源参数组合错误，提交即失败；
  \item 依赖缺失或版本冲突，反复出现；
  \item 作业排队或失败后，原因难以判断；
  \item 日志散布于各运行目录，难以定位；
  \item 报错为英文堆栈，初学者难以找到修改入口。
\end{itemize}

Slurm AI Agent 正是针对这些问题设计的。它不是简单地“让模型执行命令”，而是让模型负责它擅长的理解与解释，让程序负责它擅长的校验与执行。

\section{典型使用流程}
以下以一个典型任务为例，从上传代码到复盘结果依次经过五个环节。

\subsection{第一步：创建项目与环境准备}
用户通过 SSH 隧道打开 Web 控制台，创建项目并上传代码。每个项目对应平台上的一个目录：代码、数据放在项目根下，实验按运行子目录组织，所有提交脚本、标准输出、错误日志和报告统一归入 \code{logs/}。文件管理器支持树状浏览、在线编辑、预览和上传，路径校验拒绝越界访问项目目录之外的位置。系统在后台为项目创建专属 Conda 环境，就绪后自动通知；同项目的多个实验共享这套依赖，又各自保留独立的对话上下文。

% TODO(截图)：此处适合放文件管理器界面截图（建议文件名 agent-file-manager.png）。

\subsection{第二步：检查依赖}
点击“检查依赖”后，系统扫描项目里的依赖清单（\code{requirements.txt}、\code{environment.yml}、\code{pyproject.toml} 等）、Shell 脚本与 README 中的安装说明、以及 Python 源码中的 import 语句，生成一份依赖清单；随后向真实软件源逐项查询：该包是否存在、请求的版本是否可满足、环境是否已安装，结果直接标注在清单上。来自显式声明的可信依赖默认勾选，已安装的自动排除，拿不准的 import 候选交由模型复核并附上出处证据，供用户判断。

安装时，用户看到的不是抽象的进度条，而是 conda/pip 的真实终端输出。如果安装失败，系统先自动分类原因：求解冲突或版本不存在会自动去掉版本约束重试一次；磁盘、网络、包名错误等问题直接说明原因并终止；其余情况由模型结合真实版本查询结果给出修正方案，经用户确认后再继续。安装支持中断后增量重试，全部完成后系统逐包核验确实已安装至项目环境。

% TODO(截图)：此处适合放依赖确认清单或安装实时输出界面截图（建议文件名 agent-deps.png）。

\subsection{第三步：描述任务}
用户用自然语言描述训练或计算目标，助手读取当前目录与集群实时状态，建议资源配置并生成命令正文。项目内的对话助手接入了 19 项受控工具：查询节点、分区、QoS、作业与历史用量，检索平台文档，读取项目文件等。它回答“这个分区现在还剩多少 GPU”“我的作业为什么在排队”时，数据一律来自实时接口而非模型记忆。平台 Markdown 文档内置为本地知识库，支持语义问答；对平台概念不熟悉时，可直接向助手提问，无需在文档站中自行检索。

\subsection{第四步：确认并提交}
系统内置单卡 PyTorch、多机 DDP、CPU 批处理、Job Array、Jupyter 和通用脚本六类模板，参数默认值与范围都已预置。用户核对资源与命令后提交。提交前，后端完成一组用户容易遗漏的检查：作业名与路径是否合法、CPU 分区上是否误申请了 GPU、资源请求是否超出账户与 QoS 的授权上限、入口脚本与配置文件是否存在、运行目录是否重复。通过校验后，脚本头、运行目录、Conda 激活和日志路径由服务端锁定生成，用户的命令正文原样嵌入。

% TODO(截图)：此处适合放资源确认/提交界面截图（建议文件名 agent-submit.png）。

提交成功并不等于资源申请生效：系统会立即读取作业详情，用 TRES/GRES 证据核验 GPU 请求数，证据不足时如实标记为待确认。

\subsection{第五步：监控与复盘}
提交成功的作业进入持续观察列表，运行期间可随时查看状态、队列位置与实时日志。任务结束后，系统聚合 Slurm 状态与标准输出、错误日志的尾部，自动生成报告：成功作业给出结果概况、关键指标与警告；失败作业给出关键错误、原因分析与修复步骤。即使大模型暂不可用，报告也退化为带原始日志上下文的基础版本，保证信息不缺失。失败日志中明确的缺失依赖（如 \code{ModuleNotFoundError}）会被抽取出来写回项目依赖需求，下一次“检查依赖”时自动吸收这些运行反馈。

% TODO(截图)：此处适合放作业报告或日志详情界面截图（建议文件名 agent-report.png）。

\section{设计理念：人做决策，模型做理解，程序做约束}
系统在整体上遵循一条清晰的职责划分。

\begin{itemize}
  \item \textbf{人做决策。}安装哪些依赖、申请多少资源、何时提交、是否接受修正方案——所有有实际后果的决定都由用户确认，系统如实呈现决策所需的依据。
  \item \textbf{模型做理解。}大模型负责解释 Slurm 概念、阅读项目文件与平台文档、生成作业命令正文、诊断失败日志。它的输出是“结构化的建议”，而不是直接生效的操作。
  \item \textbf{程序做约束。}权限、配额、路径、脚本前导、环境激活与资源字段全部由后端锁定生成，模型无法绕过；提交之后系统还会回到 Slurm 读取真实结果进行核验。
\end{itemize}

另一条贯穿始终的原则是\textbf{证据优先}：队列状态、资源占用、配额这类随时变化的信息，必须通过工具实时获取，不允许模型凭记忆作答；回答保留数据来源，检索或分析失败时有确定性的降级方案，而不是沉默或编造。

\section{创新点}
以下几条均可对应到系统中的具体实现，而非泛化的工程实践。

\begin{enumerate}
  \item \textbf{模型与执行权解耦的受控提交。}大模型只产出命令正文与结构化资源字段；脚本头、运行目录、Conda 激活与日志路径由后端锁定生成；提交后再次读取作业详情，用 TRES/GRES 证据核验 GPU 请求是否生效，核验不可用时如实标记，避免“接口成功但资源未申请到”的误判。
  \item \textbf{面向具体平台的模型校准。}系统提示词内置了在该集群实测确认的平台事实：分区名称区分大小写、分区与计费账户及 QoS 的对应关系、默认配额，以及 scrontab 被禁用（周期性任务建议改为长时限作业内循环）；同时约定数据口径——当前队列数量必须由实时工具逐条统计，控制器累计统计不得混作当前状态。模型回答因此贴合本平台实际，而非通用 HPC 印象。
  \item \textbf{带证据链的知识检索。}平台 Markdown 文档按标题分块并向量化，向量缓存到磁盘；检索按余弦相似度排序，低于阈值的片段直接丢弃；回答保留来源文件与相似度，向量服务不可用时自动降级为关键词检索。
  \item \textbf{依赖安装的预检与分类自愈。}安装前向真实软件源逐项查询：包是否存在、请求版本是否可满足、环境是否已安装，针对“把其它集群 module 的版本号当成 conda 版本”的高频错误按同系列优先给出建议版本；安装失败先由确定性分类器归类：常见冲突自动去版本重试，环境类问题直接终止并说明原因，仅少数疑难情形交给模型生成修正命令，且必须经用户确认；超时采用“停滞判定”而非固定总时长，避免把大体积包的健康下载误判为卡死。
  \item \textbf{运行反馈驱动的环境闭环。}多实验共享项目依赖而会话相互隔离；作业结束后报告自动回到原会话，失败日志中确定性提取的缺失依赖写回项目需求，使下一轮“检查依赖”吸收真实运行反馈。
\end{enumerate}

\section{技术实现与安全边界}
\subsection{技术栈与规模}
项目后端采用 FastAPI，前端为无需构建的原生单页应用；聊天与安装进度使用 SSE，终端使用 WebSocket + PTY（支持 ANSI 颜色与窗口缩放，交互式调试不必另开 SSH 客户端）。源码包含 49 个 HTTP/WebSocket 路由和 19 个智能体工具。

\subsection{模型与调用方式}
\begin{itemize}
  \item 默认对话模型：\code{deepseek-v4-flash}；用户可在 Web 界面从当前 Key 可用的 deepseek/glm 模型中切换。
  \item 文档向量模型：\code{qwen3-embedding}，用于平台知识库语义检索；预留 \code{qwen3-reranker} 用于未来的检索重排序扩展。
\end{itemize}

系统通过 OpenAI Python SDK 调用学校提供的 OpenAI 兼容接口（\url{https://api.llm.ustc.edu.cn/v1}）。对话使用 Chat Completions 与 Function Calling：请求传入多轮 \code{messages} 与 19 项工具定义，模型返回 \code{tool\_calls} 时由后端执行真实工具并把结果回填给模型，直至生成最终回答。知识库使用 Embeddings 接口将 Markdown 片段与问题转换为向量，按余弦相似度检索；接口异常时自动降级为关键词检索。

\subsection{与 Slurm 的交互与安全机制}
Slurm 操作不经过大模型 API。系统通过 \code{slurmrestd v0.0.41} 调用作业、节点、QoS、历史和提交接口，使用 JWT 请求头认证；REST 未覆盖的优先级、fairshare 和实时用量信息，仅通过白名单只读 CLI 获取。在此基础上，系统实现了多项面向真实运行环境的工程机制：
\begin{itemize}
  \item 密钥来自环境变量 \code{LLM\_API\_KEY} 与 \code{SLURM\_JWT}（或由 \code{scontrol token} 自动刷新），均不写入仓库；系统默认通过 Unix Socket 与 SSH 隧道访问，不公开暴露登录节点 TCP 端口，终端连接做同源校验与并发限制；
  \item 项目级 Conda 环境采用线程锁与文件锁双层互斥，中断后的半成品环境可自愈重建；依赖安装采用停滞判定超时、失败分类器与信号阶梯终止，避免误杀健康下载或留下残缺环境；
  \item 路径穿越防护与白名单 CLI，模型与前端均无法越出项目目录或拼装任意命令；
  \item Slurm JWT 自动刷新、知识检索降级、作业报告兜底，外部服务异常时核心功能不中断。
\end{itemize}

\section{适用对象与场景}
\begin{itemize}
  \item \textbf{课程与初学者：}第一次接触 Slurm 与 Linux 集群的学生，可以在引导下完成从上传代码到拿到训练结果的完整闭环，而不必先吃透全套命令。
  \item \textbf{科研训练与批量实验：}需要反复提交、对比和复盘的研究者，可以借助项目化空间、作业模板与自动报告减少重复操作。
  \item \textbf{教学与平台管理：}系统保留了清晰、可审计的权限与执行边界，适合在多用户共享登录节点上安全部署。
\end{itemize}

对新用户，它提供贯穿整个作业生命周期的引导；对有经验的用户，它减少查看资源、组织实验和复盘日志的重复操作；对平台管理场景，它保留了清晰的权限与执行边界。系统的目标始终一致：让用户把精力集中在计算任务本身，而不是消耗在工具切换和易错配置上。

\end{document}
